% ch/maxwald3.tex: subsections of Chapter chap:maxwald (editor E3), the material of chap 570 to 895 of the chapter of N. Polson (misc/reviews/nick-maxwald/source). Notation of the merge plan, Section 2; labels mw:*; every number marked "CHECK maxwald" is replaced afterwards by the value of code/maxwald_*.py.

\section{Finite maxima and a limit argument}\label{mw:sec:finite}

The maximum interpretation of Section~\ref{mw:sec:max} suggests a direct strategy: truncate the maximum at $N$ terms, show that the finite object has no zeros off the critical line, and pass to the limit. For the sum representation instead of the maximum the book has the truncations of the Biane, Pitman and Yor lattice (Subsection~\ref{rz:sec:finite}).

\subsection{Finite maxima}

For $N\geq1$ let
\[
M_4^{(N)}=\max\Bigl\{\frac{E_n}{\pi n}\ (n\leq N),\ \ \frac{E'_m}{\pi(2m-1)}\ (m\leq N)\Bigr\},
\]
\[
\mathcal P_N(q)=\prod_{n\leq N}(1-q^n)\prod_{m\leq N}(1-q^{2m-1})=\sum_{m\geq0}p^{(N)}_mq^m,
\]
where $(E_n)_{n\geq1}$ and $(E'_m)_{m\geq1}$ are independent standard exponential variables, and $M_4=\lim_NM_4^{(N)}$ is the maximum of Subsection~\ref{sec:ch16:max}, written with the rates $\pi n$ and $\pi(2m-1)$. Then $\Pr(M_4^{(N)}\leq t)=\mathcal P_N(e^{-\pi t})$ exactly, a polynomial in $e^{-\pi t}$ with integer coefficients. Expanding the product and integrating term by term, for $\re s>0$,
\[
\E\bigl(M_4^{(N)}\bigr)^{s/2}=-\Gamma(1+\tfrac s2)\,\pi^{-s/2}D_N(s),\qquad D_N(s)=\sum_{m\geq1}p^{(N)}_m\,m^{-s/2},
\]
a finite Dirichlet polynomial $D_N$ with integer coefficients. Since the factors left out of $\mathcal P_N$ are $1-q^a$ with $a>N$, $p^{(N)}_m$ equals the coefficient of $\vartheta_4$ for $m\leq N$: it is $2(-1)^n$ at $m=n^2$ and $0$ otherwise. So the first block of $D_N$ is exactly minus the partial sum $2\sum_{n^2\leq N}(-1)^{n+1}n^{-s}$ of $2(1-2^{1-s})\zeta(s)$, followed by a tail of integer coefficients at the frequencies $\frac12\log m$ with $N<m\leq N(3N+1)/2$, the degree of $\mathcal P_N$.

By Subsection~\ref{sec:ch16:max}, $(1-2^{1-s})\zeta(s)=\E[(2K^2)^{s/2}]/(2\Gamma(1+\frac s2))$ for every complex $s$, with $\pi M_4=2K^2$ and $K$ the Kolmogorov variable, and since $2\xi(s)/(s-1)=2\pi^{-s/2}\Gamma(1+\frac s2)\zeta(s)$ this reads
\[
\E M_4^{s/2}=\frac{2(1-2^{1-s})\,\xi(s)}{s-1}\qquad\text{for every complex }s.
\]
Define the normalised finite stage and its symmetrisation
\[
\xi_N(s)=\frac{(s-1)\,\E(M_4^{(N)})^{s/2}}{2(1-2^{1-s})},\qquad \xi_N^{\mathrm{sym}}(s)=\tfrac12\bigl(\xi_N(s)+\xi_N(1-s)\bigr),
\]
so that the same expression with $M_4$ in place of $M_4^{(N)}$ is $\xi$. In the open strip $0<\re s<1$ the factor $1-2^{1-s}$ does not vanish and $\E(M_4^{(N)})^{s/2}$ is holomorphic (it is holomorphic for $\re s>-4N$, since $\Pr(M_4^{(N)}\leq t)\sim Ct^{2N}$ as $t\to0$), so $\xi_N$ and $\xi_N^{\mathrm{sym}}$ are holomorphic there. The function $\xi_N^{\mathrm{sym}}$ is invariant under $s\mapsto1-s$ and, because $\xi_N$ is real on the real axis, $\xi_N^{\mathrm{sym}}(\tfrac12+it)=\re\,\xi_N(\tfrac12+it)$ is real for real $t$.

\begin{theorem}\label{mw:thm:finite}
(i) $M_4^{(N)}\uparrow M_4$ almost surely, and $\Pr(M_4^{(N)}\neq M_4)\leq C\sqrt N\,e^{-\pi\sqrt N/2}$.

(ii) On every compact $K$ in the open strip $0<\re s<1$,
\[
\sup_K|\xi_N-\xi|\leq C_K\,N^{1/4}e^{-\pi\sqrt N/4},\qquad
\sup_K|\xi_N^{\mathrm{sym}}-\xi|\leq C_K\,N^{1/4}e^{-\pi\sqrt N/4}.
\]
\end{theorem}

\begin{proof}
(i) The variables of the stages are the same for all $N$, so adding variables can only raise the maximum, and $M_4^{(N)}\neq M_4$ requires a left out variable to exceed $M_4^{(N)}$. For any $\tau>0$,
\[
\Pr(M_4^{(N)}\neq M_4)\leq\Pr(M_4^{(N)}\leq\tau)+\sum_{n>N}e^{-\pi n\tau}+\sum_{m>N}e^{-\pi(2m-1)\tau}.
\]
The last two sums are at most $2e^{-\pi N\tau}/(1-e^{-\pi\tau})$. For the first, $\Pr(M_4^{(N)}\leq\tau)=\vartheta_4(\tau)/\Pi_N(\tau)$, where $\Pi_N(\tau)=\prod_{n>N}(1-e^{-\pi n\tau})\prod_{m>N}(1-e^{-\pi(2m-1)\tau})$ is the product of the left out factors. Since $1-x\geq e^{-x/(1-x)}$, $\Pi_N(\tau)\geq\exp\bigl(-2e^{-\pi N\tau}/((1-e^{-\pi\tau})(1-e^{-\pi N\tau}))\bigr)$, while Jacobi's transformation $\vartheta_4(\tau)=\tau^{-1/2}\sum_ne^{-\pi(n+1/2)^2/\tau}$ gives $\vartheta_4(\tau)\leq3\tau^{-1/2}e^{-\pi/(4\tau)}$ for $\tau\leq1$. The choice $\tau=1/(2\sqrt N)$ balances $\pi/(4\tau)$ against $\pi N\tau$, both equal to $\pi\sqrt N/2$, and gives the bound.

(ii) For $\sigma=\re s>0$, $M_4^{(N)}\leq M_4$ gives $|\E M_4^{s/2}-\E(M_4^{(N)})^{s/2}|\leq2\E[M_4^{\sigma/2}\mathbf 1\{M_4\neq M_4^{(N)}\}]\leq2(\E M_4^{\sigma})^{1/2}\Pr(M_4^{(N)}\neq M_4)^{1/2}$ by Cauchy and Schwarz. The prefactor $(s-1)/(2(1-2^{1-s}))$ is bounded on $K$ because $1-2^{1-s}$ vanishes only on $\re s=1$. The symmetrised bound follows since $K$ and $1-K$ both lie in the strip and $\xi(1-s)=\xi(s)$.
\end{proof}

The exponent in (ii) is not sharp. A heuristic count of the event $M_4^{(N)}\neq M_4$, in which the $j$th left out variable wins with probability $\E e^{-\pi jM_4}=\pi\sqrt j/\sinh(\pi\sqrt j)$, of order $\sqrt j\,e^{-\pi\sqrt j}$, gives a rate of order $Ne^{-\pi\sqrt N}$.

\begin{corollary}[Hurwitz]\label{mw:cor:hurwitz}
(a) Every zero of $\xi$ in the open strip is a limit of zeros of $\xi_N$, and of zeros of $\xi_N^{\mathrm{sym}}$, with multiplicity; every limit of zeros of $\xi_N$ inside the open strip is a zero of $\xi$.

(b) If $\rho=\tfrac12+i\gamma$ is a simple zero of $\xi$, then for all large $N$, $\xi_N^{\mathrm{sym}}$ has exactly one zero in a small disc about $\rho$, and that zero lies exactly on the critical line.

(c) Suppose that there are $N_j\to\infty$ and heights $h_j\to\infty$ such that $\xi^{\mathrm{sym}}_{N_j}$ has no zeros in $\{\tfrac12<\re s<1,\ |\im s|<h_j\}$ for every $j$. Then the Riemann Hypothesis holds. In particular this is so if there are $c>0$ and $N_0$ such that for every $N\geq N_0$, $\xi^{\mathrm{sym}}_N$ has no zeros in $\{\tfrac12<\re s<1,\ |\im s|<c\sqrt N\}$.
\end{corollary}

\begin{proof}
(a) is Hurwitz's theorem with Theorem~\ref{mw:thm:finite}(ii). For (b), Hurwitz gives exactly one zero $z_N$ of $\xi_N^{\mathrm{sym}}$ in a disc $D$ centred at $\rho$ for large $N$. The symmetries $s\mapsto1-s$ and $s\mapsto\bar s$ of $\xi_N^{\mathrm{sym}}$ make its zero set invariant under $s\mapsto1-\bar s$, which maps $D$ to itself; uniqueness forces $1-\bar z_N=z_N$, that is $\re z_N=\tfrac12$. For (c), fix $\Gamma$ and let $U=\{\tfrac12<\re s<1,\ |\im s|<\Gamma\}$. For $j$ so large that $h_j>\Gamma$ the functions $\xi^{\mathrm{sym}}_{N_j}$ have no zeros in $U$, and they converge to $\xi\not\equiv0$ uniformly on compact subsets of $U$, so by Hurwitz $\xi$ has no zeros in $U$. Letting $\Gamma\to\infty$ and using $\xi(s)=\xi(1-s)$ gives RH.
\end{proof}

Part (b) is pointwise in $\rho$ and gives no height up to which the stages reproduce all the zeros of $\xi$, since that would need an error bound uniform up to the height together with lower bounds for $|\xi|$ on circles around every zero below it, that is, bounds for $1/\zeta'(\rho)$, and none is proved here.

\subsection{What the finite stages look like}

The unsymmetrised stage $\xi_N$ is not zero free off the line at any of the $N$ from $4$ to $128$ that we tested, so the naive version of the strategy fails. Its zeros near $\rho_k$ approach $\rho_k$ as $N$ grows, and more slowly the larger $\gamma_k$ is (Table~\ref{mw:tab:GN}), and it carries spurious zeros elsewhere, such as the zero of $\xi_{16}$ at $1.055+17.984i$, just to the right of the strip.

\begin{table}[htbp]
\centering\small
\fitwidth{\begin{tabular}{rlll}
\toprule
$N$ & zero near $\rho_1$ & zero near $\rho_2$ & zero near $\rho_5$ \\
\midrule
16 & $0.51577+14.09915i$ & $0.95703+20.91901i$ & $1.56809+34.68128i$ \\
32 & $0.49948+14.13483i$ & $0.48159+21.02790i$ & $1.02797+32.22350i$ \\
64 & $0.50000052+14.1347256i$ & $0.50003418+21.0220191i$ & $0.51254+32.92764i$ \\
\bottomrule
\end{tabular}}
\caption{Zeros of the unsymmetrised finite stage $\xi_N$. For $N=16$ and $N=32$ the zero listed near $\rho_5$ does not track $\rho_5$, since $\rho_5$ lies above the height $4\sqrt N$ discussed below.}
\label{mw:tab:GN}
\end{table}

\begin{table}[htbp]
\centering\small
\fitwidth{\begin{tabular}{rrrrrr}
\toprule
$N$ & zeros & $|t_1-\gamma_1|$ & $|t_2-\gamma_2|$ & $|t_3-\gamma_3|$ & $|t_5-\gamma_5|$ \\
\midrule
16 & 7 & $3.6\times10^{-2}$ & $9.9\times10^{-2}$ & $1.7\times10^{-1}$ & $1.3\times10^{-3}$ \\
32 & 10 & $1.0\times10^{-4}$ & $5.9\times10^{-3}$ & $6.8\times10^{-2}$ & $1.1$ \\
64 & 8 & $4.1\times10^{-7}$ & $2.1\times10^{-5}$ & $3.8\times10^{-4}$ & $7.5\times10^{-3}$ \\
\bottomrule
\end{tabular}}
\caption{Real zeros $t_k$ of the symmetrised stage $\re\,\xi_N(\frac12+it)$: the column ``zeros'' is the number of sign changes on $[10,44]$, where $\xi$ has 8 zeros ($10\leq\gamma\leq44$). At $N=16$ the agreement at $\gamma_5$ is accidental.}
\label{mw:tab:sym}
\end{table}

From $N=32$ to $N=64$ the error at $\gamma_1$ falls by a factor of about $250$, against $e^{\pi(8-\sqrt{32})}=1574$ for the exponential factor of the heuristic rate (about $800$ with its prefactor $N$) and $e^{\pi(8-\sqrt{32})/4}=6.3$ for that of the proven one; this is a rate of about $e^{-0.75\pi\sqrt N}$ (about $e^{-0.60\pi\sqrt N}$ between $N=8$ and $N=16$). The supremum of $|\xi_N-\xi|$ over the grid $\{0.3,0.5,0.7\}+i\{0,2,\dots,10\}$ decays like $e^{-c\pi\sqrt N}$ with $c\approx0.93$ between $N=32$ and $N=64$ and $c\approx0.80$ between $N=8$ and $N=16$, and that of $|\xi_N^{\mathrm{sym}}-\xi|$ with $c\approx0.81$ and $c\approx0.80$, all between the exponent $\frac14$ of the proof and the exponent $1$ of the heuristic. Between consecutive stages the exponents for $\xi_N^{\mathrm{sym}}$ and for the error at $\gamma_1$ oscillate, between $0.73$ and $1.10$ and between $0.49$ and $1.37$, so the exponent $1$ of the heuristic is not a bound. Since $|\xi(\tfrac12+i\gamma)|$ decays like $e^{-\pi\gamma/4}$, the heuristic predicts that zeros up to height $\gamma$ are captured once $\pi\sqrt N$ exceeds $\pi\gamma/4$ by a margin, that is $N\gtrsim\gamma^2/16$. The tables agree: the symmetrised stage $N=64$ captures every zero up to height $33$ to within $10^{-2}$ (and $\xi_{64}$ itself to within $1.5\times10^{-2}$), and the accuracy degrades beyond $37$. This suggests that the stage of index $N$ tracks the zeros of $\xi$ up to a height of about $4\sqrt N$; this tracking height is not the zero free height of Corollary~\ref{mw:cor:hurwitz}(c), which may grow with $N$ in any way.

\subsection{The symmetric stages $F_N^{\mathrm{sym}}$ and how far they stay on the line}\label{mw:sec:finiteN}

The normalisation $\xi_N$ divides by $1-2^{1-s}$, which has zeros on $\re s=1$ and, after symmetrisation, on $\re s=0$. For counting zeros it is better to keep the factor. Put
\[
F_N(s)=\frac{s-1}2\,\E\bigl(M_4^{(N)}\bigr)^{s/2}=-\frac{s-1}2\,\Gamma\bigl(1+\tfrac s2\bigr)\pi^{-s/2}D_N(s),
\]
\[
F_N^{\mathrm{sym}}(s)=\tfrac12\bigl(F_N(s)+F_N(1-s)\bigr),
\]
so that $F_N=(1-2^{1-s})\xi_N$. These functions are holomorphic on $-4N<\re s<4N+1$, but they are not entire: $F_N$ has a pole at $s=-4N$ and $F_N^{\mathrm{sym}}$ also at $s=4N+1$, far from every region used below. As $N\to\infty$, $F_N\to(1-2^{1-s})\xi(s)$ and $F_N^{\mathrm{sym}}\to\mathfrak q(s)\xi(s)$, with $\mathfrak q(s)=1-2^{-s}-2^{s-1}$. The zeros of $\mathfrak q$ are simple and are the solutions of $2^s=1\pm i$, all on $\re s=\tfrac12$, at $t=(2\pi k\pm\pi/4)/\log2$.

Theorem~\ref{mw:thm:finite} and Corollary~\ref{mw:cor:hurwitz} hold for $F_N^{\mathrm{sym}}$ with $\mathfrak q\xi$ in place of $\xi$, on compact subsets of the open strip, since $1-2^{1-s}$ is bounded there. Part (c) becomes: if there are $N_j\to\infty$ and $h_j\to\infty$ such that $F^{\mathrm{sym}}_{N_j}$ has no zeros in $\{\frac12<\re s<1,\ |\im s|<h_j\}$, then RH holds, by the same argument, since a zero of $\xi$ in that region is a zero of $\mathfrak q\xi$. We call (c) for $F_N^{\mathrm{sym}}$ the $F_N^{\mathrm{sym}}$ form of Corollary~\ref{mw:cor:hurwitz}(c): it is a different statement from (c) for $\xi_N^{\mathrm{sym}}$, since $F_N^{\mathrm{sym}}$ is not $\xi_N^{\mathrm{sym}}$ times a fixed factor.

Zeros of $F_N^{\mathrm{sym}}$ in a rectangle $[-1,2]\times[t_1,t_2]$ can be counted by the argument principle, and zeros on the line are sign changes of the real function $F_N^{\mathrm{sym}}(\tfrac12+it)$. When the two counts agree, every zero in the rectangle is on the line. Write $T^*(N)$ for the height below which the counts agree in every window of height $10$ and above which off line zeros appear in the rectangle. Since the rectangle contains the strip, $F_N^{\mathrm{sym}}$ has no zeros in $\{\tfrac12<\re s<1,\ |\im s|<T^*(N)\}$, which is the hypothesis of the $F_N^{\mathrm{sym}}$ form of Corollary~\ref{mw:cor:hurwitz}(c) at the height $T^*(N)$. Table~\ref{mw:tab:Tstar} reports the counts.

\begin{table}[htbp]
\centering\small
\fitwidth{\begin{tabular}{rrrrrr}
\toprule
$N$ & terms & height & on the line & first off line window & $T^*(N)$ \\
\midrule
8 & 72 & 400 & 70 of 192 & $(110,120]$ & $\approx110$ \\
16 & 350 & 1000 & 146 of 510 & $(250,260]$ & $\approx250$ \\
32 & 1478 & 800 & 312 of 360 & $(530,540]$ & $\approx530$ \\
48 & 3372 & 1000 & 494 of 496 & $(870,880]$ & $\approx870$ \\
64 & 6034 & 1500 & 791 of 797 & $(1060,1070]$ & $\approx1060$ \\
96 & 13670 & 2500 & 1463 of 1465 & $(1590,1600]$ & $\approx1590$ \\
128 & 24354 & 130 & 64 of 66 & $(60,70]$ & $\approx60$ \\
\bottomrule
\end{tabular}}
\caption{Zeros of the symmetric stage $F_N^{\mathrm{sym}}$ in $[-1,2]\times[0,T]$ by the argument principle, against sign changes on the critical line: ``terms'' is the number of nonzero coefficients $p^{(N)}_m$, $m\geq1$, ``height'' the height $T$ searched, and ``on the line'' the number of sign changes out of the number of zeros counted. The windows have height $10$. Above $T^*(N)$ the off line zeros arrive in bulk for $N\leq32$: for $N=16$ the two windows of $(300,320]$ have 10 of their 13 zeros off the line. The sum $D_N$, whose terms are much larger than the sum itself, is evaluated in double precision, with a bound on its rounding error on the critical line in every window, and in multiprecision in the windows below $100$ for $N=96$ and below $130$ for $N=128$; above $130$ the double precision margin of $N=128$ is lost, so that stage is reported only up to $130$. Off line zeros come in pairs $\sigma+it$, $1-\sigma+it$, so an off line pair adds two to the count of its window, and a difference of one between the two counts is an artefact, such as a zero at the window boundary.}
\label{mw:tab:Tstar}
\end{table}


The first off line pairs of $F_N^{\mathrm{sym}}$ are $1.5643+111.7906i$ for $N=8$, $1.1950+254.3196i$ for $N=16$, $1.2079+537.2287i$ for $N=32$, $1.0434+870.3807i$ for $N=48$, $1.1025+1063.9243i$ for $N=64$, $1.1177+1593.1234i$ for $N=96$ and $1.5311+60.7314i$ for $N=128$, each with its mirror image under $s\mapsto1-\bar s$. The first off line zero inside the strip appears for $N=8$ in $(120,130]$, for $N=16$ in $(330,340]$ and for $N=32$ in $(550,560]$, and none was found there for $N=48$, $64$, $96$ and $128$ up to the heights $1000$, $1500$, $2500$ and $130$ searched, so the height up to which $F_N^{\mathrm{sym}}$ has no off line zeros in the strip is at least $T^*(N)$ and may be larger. The counts depend on the rectangle: for $N=16$ the rectangle $[-8,9]\times[240,250]$ contains $7$ zeros where $[-1,2]\times[240,250]$ contains $3$, all on the line, so zeros of $F_N^{\mathrm{sym}}$ far from the strip exist below $T^*(N)$, and $T^*(N)$ refers to $[-1,2]$.

Three things follow. First, the symmetric stages are not real rooted: off line zeros exist in $[-1,2]\times[0,T]$ for every $N$ we tested, but only above a height $T^*(N)$, and the first ones found lie outside the critical strip. Second, up to $N=96$ the height $T^*(N)$ grows roughly linearly, while the height up to which the stage approximates $\mathfrak q\xi$ well enough to capture its zeros is only about $4\sqrt N$ (Table~\ref{mw:tab:sym} is for $\xi_N^{\mathrm{sym}}$; the same height is expected for $F_N^{\mathrm{sym}}$, whose limit $\mathfrak q\xi$ has the same zeros besides those of $\mathfrak q$). At $N=128$ the growth breaks: $T^*(128)/128=0.47$, although the strip has no off line zero of that stage below $130$, the height to which it was counted. Below $4\sqrt N$ the zeros of $F^{\mathrm{sym}}_N$ are the zeros of $\mathfrak q\xi$ moved along the line; between $4\sqrt N$ and $T^*(N)$ the stage has its own zeros, fewer than $\mathfrak q\xi$ has, but still all on the line; above $T^*(N)$ it has zeros off the line. The linear growth is what a heuristic based on the approximate functional equation suggests for a symmetrised Dirichlet polynomial whose longest frequency is $\frac12\log m_{\max}$ with $m_{\max}=N(3N+1)/2$: the two halves of $F_N^{\mathrm{sym}}$ can balance off the line only when $(t/2\pi)^{\sigma-1/2}$ is matched by a ratio of the Dirichlet polynomial at $\sigma$ and $1-\sigma$, which first becomes possible at $t$ of order $N$; the pair of $N=128$, far from the strip, falls outside this heuristic. Third, the hypothesis of the $F_N^{\mathrm{sym}}$ form of Corollary~\ref{mw:cor:hurwitz}(c) holds up to height $T^*(N)$ for the $N$ tested, but finitely many $N$ cannot verify an asymptotic hypothesis, and (c) needs no particular height: it suffices that $T^*(N_j)\to\infty$ along some sequence $N_j\to\infty$. The finite $N$ question becomes the growth of $T^*(N)$: the conjecture $T^*(N)\geq\kappa N$ for some $\kappa>0$ and all $N$ would imply RH, as would any $T^*(N_j)\to\infty$, in particular the far weaker $T^*(N)\geq c\sqrt N$.

We do not have a mechanism for $T^*(N)$. The classical argument that puts zeros of a symmetrised Dirichlet polynomial on the line, dominance of the first term of $D_N$, that is $|D_N(s)-p_1|<|p_1|$ on $\re s\geq\tfrac12$ (with $|p_1|=2$), fails here because $\sum_m|p^{(N)}_m|m^{-1/4}$ is large: about $50$ for $N=8$ and $2.4\times10^9$ for $N=64$. The observed $T^*(N)$ is nevertheless what the balance heuristic predicts, so the right statement to prove may be about the ratio $|D_N(\sigma+it)/D_N(1-\sigma+it)|$ rather than about $D_N$ alone.

\subsection{Comparison with the symmetrised lattice truncations}

Subsection~\ref{rz:sec:finite} carries out the same strategy for the symmetrised lattice truncations $K_n$, with Proposition~\ref{rz:prop:symtrunc} and Theorem~\ref{rz:thm:symhurwitz} as the analogues of Corollary~\ref{mw:cor:hurwitz}(b) and (c); the off line zeros of $F_N^{\mathrm{sym}}$ above $T^*(N)$ do not bear on Conjecture~S of that subsection, which concerns $K_n$.

\section{Bondesson classes containing GGC}\label{mw:sec:bondesson}

The classes BO, GGC and SD (Subsection~\ref{sec:ch3:hierarchy}) have finite order versions. This section establishes, for the clock $T$, the finite order versions up to the orders $1.16\times10^{25}$ on the L\'evy side and $9\times10^{12}$ on the transform side from finite verification, and shows that every infinite order version is equivalent to RH.

\begin{definition}\label{mw:def:orders}
A function $f$ on an interval is $k$ monotone if $(-1)^jf^{(j)}\geq0$ for $0\leq j\leq k$, and completely monotone if it is $k$ monotone for every $k$. For $k\geq1$ let $\mathrm{GGC}_k$ be the laws on $[0,\infty)$ with a L\'evy density $\ell$ such that $x\ell(x)$ is $k$ monotone, and $\mathrm{BO}_k$ those with $\ell$ itself $k$ monotone. Thus $\mathrm{GGC}_1=\mathrm{SD}$ is the self decomposable class, $\bigcap_k\mathrm{GGC}_k=\mathrm{GGC}$ and $\bigcap_k\mathrm{BO}_k=\mathrm{BO}$. A positive function $\varphi$ on $(0,\infty)$ is HCM of order $k$ if for every $u>0$ the function $w\mapsto\varphi(uv)\varphi(u/v)$, $w=v+1/v\geq2$, is $k$ monotone on $[2,\infty)$; HCM means HCM of every order (Subsection~\ref{sec:weak:hcm}).
\end{definition}

The finite order versions are introduced here. They are not Urbanik's nested classes, and other finite order notions in the literature (Williamson's $k$ monotone functions, Bondesson's classes of finite order for densities) use other conventions. By the Leibniz identity $(-1)^j(fg)^{(j)}=\sum_i\binom ji[(-1)^if^{(i)}][(-1)^{j-i}g^{(j-i)}]$ a product of two nonnegative $k$ monotone functions is $k$ monotone; as $x^{-1}$ is completely monotone and $\ell(x)=x\ell(x)\cdot x^{-1}$, this gives $\mathrm{GGC}_k\subset\mathrm{BO}_k$ for every $k$.

\begin{theorem}\label{mw:thm:bondesson}
Let $T$ be the clock, with L\'evy density $W(x)/x$ and Laplace transform $\varphi_{1/2}(\lambda)=\xi(\tfrac12)/\xi(\tfrac12+\sqrt\lambda)=1/G(\lambda)$, and let $H=3\times10^{12}$. Statements (i) and (ii) are unconditional: they use the verification of RH to height $H$ \cite{PT21}.
\begin{enumerate}
\item[(i)] $T\in\mathrm{GGC}_k$, and hence $T\in\mathrm{BO}_k$, for every $k\leq1.16\times10^{25}$.
\item[(ii)] $\varphi_{1/2}$ is HCM of order $k$ for every $k\leq9\times10^{12}$.
\item[(iii)] The following are equivalent: RH; $T$ is a GGC; $T\in\mathrm{BO}$; $\varphi_{1/2}$ is HCM; $\varphi_{1/2}$ is HCM of every order.
\end{enumerate}
\end{theorem}

\begin{proof}
(i) By Theorem~\ref{thm:ch12:idsd} the clock $T$ exists and is infinitely divisible with L\'evy density $\ell(x)=W(x)/x$, and Theorem~\ref{thm:ch12:window} gives $(-1)^jW^{(j)}(x)>0$ for every $x>0$ and every integer $0\leq j\leq1.29(H-23)^2$, which exceeds $1.16\times10^{25}$. Since $x\ell(x)=W(x)$, this says $T\in\mathrm{GGC}_k$ for every $k\leq1.16\times10^{25}$, and the Leibniz step above gives $T\in\mathrm{BO}_k$.

(ii) Zeros are counted with multiplicity, and $\rho=\tfrac12+b+i\gamma$ with $\gamma>0$ runs over the zeros of positive ordinate. Put $a_\rho=-(\rho-\tfrac12)^2=\gamma^2-b^2-2ib\gamma$, as in Section~\ref{ch:clock}, so that $\re a_\rho>0$ and $\sum_\rho|a_\rho|^{-1}<\infty$, and by the Hadamard product \eqref{eq:ch2:hadamard}, $\varphi_{1/2}(\lambda)=\prod_\rho(1+\lambda/a_\rho)^{-1}$, absolutely convergent. Since $(1+uv/a)(1+u/(va))=1+uw/a+u^2/a^2$ with $w=v+1/v$,
\[
\begin{aligned}
\varphi_{1/2}(uv)\,\varphi_{1/2}(u/v)&=\prod_{b=0,\ \gamma>0}\Bigl(1+\frac{uw}{a_\rho}+\frac{u^2}{a_\rho^2}\Bigr)^{-1}\\
&\quad\times\prod_{b>0,\ \gamma>0}\Bigl|1+\frac{uw}{a_\rho}+\frac{u^2}{a_\rho^2}\Bigr|^{-2},\qquad w=v+\frac1v,
\end{aligned}
\]
where in the second product $\rho$ stands for the pair $\{\rho,1-\bar\rho\}$ of zeros with $\gamma>0$ and $b\neq0$, whose values $a_\rho$ and $a_{1-\bar\rho}=\bar a_\rho$ are conjugate, so that their two factors multiply to a squared modulus. Put $c_\rho(u)=a_\rho/u+u/a_\rho$, so that $1+uw/a+u^2/a^2=(u/a)(w+c)$. For a zero on the line $a_\rho=\gamma^2$ and $c_\rho(u)\geq2$ is real, and the factor $(a_\rho/u)/(w+c_\rho(u))$ is completely monotone in $w$ on $[2,\infty)$. For a pair the factor is $(|a_\rho|^2/u^2)/((w+c)(w+\bar c))$ with $c=c_\rho(u)$ complex. Write $a_\rho=re^{-i\phi}$ with $0<\phi<\pi/2$. Then $\tan\phi=2b\gamma/(\gamma^2-b^2)\leq\gamma/(\gamma^2-\tfrac14)<1.01/\gamma$ for $\gamma\geq10$, because $0<b<\tfrac12$. Moreover $\re c=(r/u+u/r)\cos\phi$ and $|\im c|=|r/u-u/r|\sin\phi\leq(r/u+u/r)\sin\phi$, so for every $w\geq0$ the angle $\vartheta=|\arg(w+c)|$ satisfies $\tan\vartheta=|\im c|/(w+\re c)\leq\tan\phi$, that is $\vartheta\leq\phi$. Partial fractions give $1/((w+c)(w+\bar c))=-\im[(w+c)^{-1}]/\im c$, whence
\[
(-1)^j\frac{d^j}{dw^j}\frac1{(w+c)(w+\bar c)}=\frac{j!\,\sin((j+1)\vartheta)}{|\im c|\,|w+c|^{j+1}},
\]
which is positive as long as $(j+1)\vartheta<\pi$, hence for all $j<\pi/\phi-1$ (when $\im c=0$, that is $u=r$, the factor is $(w+c)^{-2}$, whose signed derivatives $(j+1)!\,(w+c)^{-j-2}$ are positive). By the verification of RH to height $H$ \cite{PT21} every zero with $\gamma\leq H$ is on the line, so every pair has $\gamma>H$ and $\phi<1.01/H$, and its factor is $k$ monotone for every $k\leq\pi H/1.01-1$, which is $9.33\times10^{12}$ and exceeds $9\times10^{12}$. Products of nonnegative $k$ monotone functions are $k$ monotone by the Leibniz identity, and the partial products converge, with all their $w$ derivatives, locally uniformly on $\re w>0$: the factors have no zeros there, since the poles $w=-c$ lie in $\re w<0$, and $\sum_\rho|uw/a_\rho+u^2/a_\rho^2|<\infty$ locally uniformly. So the limit is $k$ monotone.

(iii) RH holds if and only if $T$ is a GGC by Proposition~\ref{prop:rigid}, and if and only if $\varphi_{1/2}$ is HCM by Corollary~\ref{cor:hcm}(ii), that is by Bondesson's theorem \cite[Theorem 6.1.1]{Bondesson} (a function $\varphi$ with $\varphi(0+)=1$ is the Laplace transform of a GGC if and only if it is HCM; see also \cite{JRY}). A GGC is in BO by the Leibniz step. Conversely suppose that $\ell(x)=W(x)/x$ is completely monotone, $\ell(x)=\int_{[0,\infty)}e^{-xy}m(dy)$ with $m\geq0$. Taking Laplace transforms of $W(x)=x\ell(x)$ for $\re\lambda>0$,
\[
\int_0^\infty e^{-\lambda x}W(x)\,dx=\int_{[0,\infty)}\frac{m(dy)}{(\lambda+y)^2}=\sum_\rho\frac1{\lambda+a_\rho}=\frac{G'}G(\lambda).
\]
The middle expression is analytic on $\mathbb C\setminus(-\infty,0]$, while $G'/G$ has a pole at every $-a_\rho$, and $-a_\rho\notin(-\infty,0]$ exactly when $a_\rho\notin(0,\infty)$, that is when $\rho$ is off the line. So BO forces RH.
\end{proof}

For a model pair factor with a zero at $\beta=0.6$, $\gamma=100$, the first order at which the closed form fails is $1571$ at $u=1$, $w=2$, and $1570$ as $u\to0$, against the bound $\pi\gamma/(2(\beta-\frac12))=1570.8$.

\section{Is there a new path to the Riemann Hypothesis?}\label{mw:sec:path}

The honest answer is no.

\subsection{Three reformulations and what each costs}

The derivative ladder: RH holds if and only if $W$ is completely monotone (Theorem~\ref{mw:thm:stieltjes}(e)), and order $k$ and height $H$ trade quadratically (Subsection~\ref{sec:ch12:reach}), so climbing the ladder to infinity is verifying RH to infinite height; the ladder is a reformulation, not a path. The basepoint family: the basepoint clock $T_\alpha$ (Section~\ref{mw:sec:basepoint}) is GGC for every $\alpha>\tfrac12$ if and only if RH (Theorem~\ref{ri:thm:ggcriesz}), and this is Lagarias's criterion \cite{Lagarias} (Remark~\ref{rem:ch2:oneline}), so the Thorin route at basepoints and the Lagarias route are one route. The finite maxima: Corollary~\ref{mw:cor:hurwitz}(c) reduces RH to zero freeness of the symmetrised stages $\xi_N^{\mathrm{sym}}$ off the line up to heights $h_j\to\infty$, for infinitely many $N$, a hypothesis that is not weaker than RH (whether RH implies it is open). Unsymmetrised sections of $\zeta$ are known to have zeros off the line and even to the right of $\re s=1$ \cite{Turan48,Montgomery83}, which is why the symmetrisation is essential; for the symmetrised sections there are results of Spira and of Gonek and Montgomery \cite{Spira67,GonekMontgomery13} (the paragraph after Table~\ref{rz:tab:symwin}); we know of no result of this kind for the symmetrised stages of the maximum.

\subsection{The arithmetic form below \texorpdfstring{$\alpha=1$}{alpha = 1}, and how far it sees}

The sine squared prime form of Section~\ref{ch:clock} (the paragraph after \eqref{eq:sin2}, with Table~\ref{tab:prime}) writes $v_\alpha(\theta)=\re\,\xi'/\xi(\alpha+i\theta)$ for $\alpha>1$ through the signed canonical measure with atoms $\Lambda(n)n^{-\alpha}$ at $\log n$, an archimedean density and a pole density: the prime and gamma terms are nonnegative and the pole terms negative, so the positivity of $v_\alpha$ for $\alpha>1$, which is Lagarias positivity there, comes from the absence of zeros with $\beta\ge1$ and not from the terms one by one. Theorem~\ref{ri:thm:ggcriesz} already continues it below $\alpha=1$ by compensation: its last condition, that the compensated sine squared sum converges for every $\alpha'>\alpha$, is equivalent to zero freeness above $\alpha$. A smoothed form of the continuation, with explicit pole terms, is the following.

\begin{proposition}\label{mw:prop:AX}
Let $\alpha=\tfrac12+\delta$ with $\delta>0$, $t\in\R$ and $X>0$. Put $g_X(x)=e^{-\delta|x|}\cos(tx)\,(1-|x|/X)_+$ and $h_X(r)=\int_\R g_X(x)e^{irx}dx$, so that $h_X\geq0$ on $\R$: it is $(2\pi)^{-1}$ times the convolution of the two sided Poisson kernel $h_\infty(r)=\delta/(\delta^2+(t-r)^2)+\delta/(\delta^2+(t+r)^2)$ with the Fej\'er kernel $4\sin^2(Xr/2)/(Xr^2)$. Define
\[
\begin{aligned}
A_X(\delta,t)=\frac12\Bigl[&h_X(\tfrac i2)+h_X(-\tfrac i2)-\log\pi+\frac1{2\pi}\int_\R h_X(r)\,\re\psi\bigl(\tfrac14+\tfrac{ir}2\bigr)\,dr\\
&-2\sum_{n<e^X}\frac{\Lambda(n)}{n^{\alpha}}\cos(t\log n)\Bigl(1-\frac{\log n}X\Bigr)\Bigr],
\end{aligned}
\]
where $\psi$ is the digamma function and $h_X(\tfrac i2)+h_X(-\tfrac i2)=2\re[J(-\delta+\tfrac12+it)+J(-\delta-\tfrac12+it)]$ with $J(c)=-1/c+(e^{cX}-1)/(Xc^2)$. Then:
\begin{enumerate}
\item[(a)] $A_X(\delta,t)=\frac12\sum_\rho h_X(\gamma_\rho)$, with $\gamma_\rho=(\rho-\tfrac12)/i$ and the sum over all nontrivial zeros, unconditionally.
\item[(b)] Under RH, $A_X(\delta,t)\geq0$ for every $X$, $\delta$ and $t$. In particular a single triple $(\delta,t,X)$ with $A_X(\delta,t)<0$ would disprove RH.
\item[(c)] If every zero has $|\beta-\tfrac12|<\delta$ (under RH, or unconditionally when $\delta\geq\tfrac12$), then $A_X(\delta,t)\to\re\,\xi'/\xi(\alpha+it)$ as $X\to\infty$.
\end{enumerate}
\end{proposition}

\begin{proof}
(a) This is Weil's explicit formula $\sum_\rho h(\gamma_\rho)=h(\tfrac i2)+h(-\tfrac i2)-g(0)\log\pi+\frac1{2\pi}\int_\R h(r)\re\psi(\tfrac14+\tfrac{ir}2)\,dr-2\sum_n\Lambda(n)n^{-1/2}g(\log n)$, in the form of \cite{IwaniecKowalski04}, which evaluates the distribution $D$ of Proposition~\ref{rz:thm:weil} on a test function $g$ with Fourier transform $h$, applied to the even test function $g_X$, which is continuous, compactly supported and piecewise smooth, with Fourier transform $h_X$, so that $g_X(0)=1$ and $n^{-1/2}g_X(\log n)=n^{-\alpha}\cos(t\log n)(1-\log n/X)_+$; the sum over zeros converges absolutely because $h_X(r)=O(r^{-2})$ uniformly in $|\im r|\leq\tfrac12$. Here $J(c)=\int_0^Xe^{cx}(1-x/X)\,dx$, and $h_X(r)=2\int_0^Xe^{-\delta x}\cos(tx)\cos(rx)(1-x/X)\,dx=\tfrac12\sum_{s_1,s_2=\pm1}J(-\delta+i(s_1t+s_2r))$. At $r=\pm\tfrac i2$ two of the four terms are the complex conjugates of the other two, so $h_X(\pm\tfrac i2)=\re[J(-\delta+\tfrac12+it)+J(-\delta-\tfrac12+it)]$ and the two pole terms add up to twice that, which is the stated value.

(b) Under RH every $\gamma_\rho$ is real and $h_X\geq0$ on $\R$, so every term of (a) is nonnegative.

(c) For a zero $\rho=\tfrac12+b+i\gamma$ the four exponentials of $h_X(\gamma_\rho)=h_X(\gamma-ib)$ have $\re c=-\delta\pm b$, which is negative when $|b|<\delta$. Then $J(c)\to-1/c$ as $X\to\infty$, and $|J(c)+1/c|\leq2/(X|c|^2)$, so $h_X(\gamma_\rho)\to h_\infty(\gamma_\rho)=\delta/(\delta^2+(t-\gamma_\rho)^2)+\delta/(\delta^2+(t+\gamma_\rho)^2)$, with $|h_X(\gamma_\rho)|\leq C(1+|\gamma_\rho|)^{-2}$ uniformly in $X\geq1$, where $C$ depends on $\delta$, $t$ and the finitely many zeros with $||\gamma_\rho|-t|\leq1$, because the four terms $-1/c$ combine in pairs to $O(\gamma^{-2})$ and $|c|\geq1$ for the other zeros. By dominated convergence $A_X\to\frac12\sum_\rho h_\infty(\gamma_\rho)$. Finally $h_\infty(\gamma_\rho)$ is the sum of the four terms $\frac12/(\alpha\pm it-\rho)$ and $\frac12/(\rho-(1-\alpha)\pm it)$, and since the zeros are symmetric under $\rho\mapsto1-\rho$ and $\rho\mapsto\bar\rho$, $\frac12\sum_\rho h_\infty(\gamma_\rho)=\re\sum_\rho(\alpha+it-\rho)^{-1}=\re\,\xi'/\xi(\alpha+it)$ by the symmetric Hadamard product.
\end{proof}

\begin{remark}\label{mw:rem:Xhypothesis}
The hypothesis of (c) is needed. If a zero has $b=\beta-\tfrac12>\delta$, the terms of $h_X(\gamma-ib)$ with $\re c=b-\delta>0$ contain $e^{(b-\delta)X}/(Xc^2)$, which grows with $X$, and $h_\infty(\gamma-ib)$ is only the analytic continuation of a divergent integral. At $t=\gamma$ the four zeros $\tfrac12\pm b\pm i\gamma$ contribute $e^{(b-\delta)X}/(X(b-\delta)^2)\,(1+O((b-\delta)^2/\gamma^2))+O(1)$ to $A_X(\delta,\gamma)$, which is positive and unbounded in $X$. No converse is claimed: whether an off line zero forces $A_X(\delta,t)<0$ somewhere for large $X$ is not shown here.
\end{remark}

The proposition is Lagarias's criterion restricted to a family, and it is also a form of Weil's positivity criterion (compare Proposition~\ref{rz:thm:weil}) for the positive definite test functions $g_X$, which have corners. Table~\ref{mw:tab:AX} gives the minimum over $t\in[0,60]$, positive at every $X$ and every $\alpha$ tested.

\begin{table}[htbp]
\centering\small
\begin{tabular}{rrrr}
\toprule
$X$ & $\alpha=0.75$ & $\alpha=0.6$ & $\alpha=0.55$ \\
\midrule
2 & 0.03275 & 0.02495 & 0.02230 \\
4 & 0.02208 & 0.01432 & 0.01160 \\
8 & 0.01707 & 0.00954 & 0.00681 \\
13 & 0.01510 & 0.00812 & 0.00576 \\
\bottomrule
\end{tabular}
\caption{$\min_{0\leq t\leq60}A_X(\alpha-\tfrac12,t)$ from primes below $e^X$ (up to $4.4\times10^5$ at $X=13$). The values decrease towards $\xi'/\xi(\alpha)=0.011551$, $0.0046209$, $0.0023105$ for the three $\alpha$.}
\label{mw:tab:AX}
\end{table}

The positivity in the table is not evidence of anything new, and it is important to say why. The contribution of zeros above height $H$ to $A_X$ is at most about $(e^{X/2}/X+\delta)\log H/H$, because $|h_X(\gamma_\rho)|\lesssim e^{|\im\gamma_\rho|X}/(X\gamma^2)$ and $|\im\gamma_\rho|\leq\tfrac12$. With $H=3\times10^{12}$ and $X=13$ that is below $10^{-8}$, about $4.9\times10^{-10}$, far under the minima in the table. So the verified zeros force $A_X\geq0$ for $t$ in the range of the table as long as this bound stays below the minima, that is up to $X$ of about $46$ to $50$, where it reaches the values from $0.0058$ down to the limit $0.0023$; the bound itself stays small up to about $2\log H\approx57$. The arithmetic test at length $e^X$ sees zeros up to a height of order $e^{X/2}$, the usual duality between primes and zeros, and to see past the present verification one would need primes up to about $10^{25}$. As a computation the arithmetic form is strictly weaker than zero verification. As a target for proof it is sharper than before: RH needs $A_X\geq0$ uniformly in $X$ (whether this suffices is open, Remark~\ref{mw:rem:Xhypothesis}), with the prime sum, the gamma integral and the pole terms each explicit, and the whole difficulty is that the pole terms, of size $e^{(1/2-\delta)X}/X$, are cancelled by the prime sum only through the prime number theorem with an error term, which is to say through the zeros.

\section{Open problems}\label{mw:sec:open}

\begin{problem}\label{mw:prob:hurwitz}
Prove the hypothesis of the $F_N^{\mathrm{sym}}$ form of Corollary~\ref{mw:cor:hurwitz}(c): find heights $h_N\to\infty$ such that for infinitely many $N$ the symmetric stage $F^{\mathrm{sym}}_N$ has no zeros with $\tfrac12<\re s<1$ and $|\im s|<h_N$, for instance $h_N=c\sqrt N$ with some $c>0$. Subsection~\ref{mw:sec:finiteN} finds, for the $N$ tested, that every zero of the stages with $-1\le\re s\le2$ and $0<\im s\le T^*(N)$ lies on the line, with $T^*(N)$ of order $N$ for $N\le96$, far more than needed (at $N=128$ a pair outside the strip comes down to height $60.7$, while the strip has none below $130$), and the sharper conjecture is $T^*(N)\geq\kappa N$. Either statement is about explicit functions built from the Dirichlet polynomials $D_N$ with integer coefficients, and either implies RH. A mechanism for $T^*(N)$, presumably through the ratio $|D_N(\sigma+it)/D_N(1-\sigma+it)|$, is the first thing to find. The corresponding statement for the symmetrised lattice truncations $K_n$ is Conjecture~S of Subsection~\ref{rz:sec:finite} (Theorem~\ref{rz:thm:symhurwitz}).
\end{problem}

\begin{problem}\label{mw:prob:lagarias}
Theorem~\ref{ri:thm:ggcriesz} continues the sine squared prime form below $\alpha=1$ by compensation, and shows that the compensated sum converges for every $\alpha'>\alpha$ exactly when $\zeta$ has no zero with $\re s>\alpha$; Proposition~\ref{mw:prop:AX} gives a smoothed form of the continuation with explicit pole terms, which is not positive term by term. Is there a continuation of the prime form below $\alpha=1$ that is positive term by term? The prime sum diverges there, but smoothed versions of it, with $n^{-\alpha}$ replaced by a weight that cuts off at $n\approx e^{1/\varepsilon}$, remain positive termwise, and the question is whether the loss from the cutoff can be controlled by $(\xi'/\xi)(\alpha)$, the gamma sum and the pole terms. A positive answer on $\alpha\in(\alpha_0,1]$ would prove the zero free region $\beta\leq\alpha_0$.
\end{problem}

\begin{problem}\label{mw:prob:flow}
For $0<t<0.2$, does the law $T_t$ of Section~\ref{mw:sec:weak}, half of the flow clock of Section~\ref{ch:dbn}, with transform $R_t(\lambda/2)$, exist as a probability law, and is it infinitely divisible? Unconditionally it is known to be a GGC only for $t\geq0.2$, and to exist only for $t\geq0.2$ and at $t=0$, where $T_0=T/2$ (Theorem~\ref{thm:ch12:idsd}); under RH it exists and is a GGC for every $t\geq0$. This is the analogue along the de Bruijn Newman flow of the existence of $T$ (Theorem~\ref{thm:ch12:idsd}); it would need an explicit formula for the heat trace of the flow, which the Gaussian tilt destroys.
\end{problem}

\begin{remark}\label{mw:rem:flowcert}
No published result gives an explicit height below which the zeros of $H_t$ are real for $0<t<0.2$: \cite{Polymath} and \cite{PT21} give $\Lambda\leq0.2$ through a barrier at $|\operatorname{Im}z|\geq y_0>0$, and \cite{KKL09} show that all but finitely many zeros are real, without a height. Real zeros do not collide forward in time (a multiple real zero splits into real zeros, as in Lemma~\ref{ft:lem:split}, after the statement of \cite{deBruijn} cited in Section~\ref{ch:dbn}), so the verification of RH to $3\times10^{12}$ passes to $H_t$ except for zeros entering from above, which a box barrier touching the real axis would exclude. That route has not been carried out. The forward statement is Lemma~\ref{ft:lem:split}, and the barrier hypothesis and its consequence are in Proposition~\ref{ft:prop:below}. With a certified height $H'$, an explicit count of the zeros of $H_t$ and the low zeros of $H_t$ located, the certificate of Section~\ref{sec:ps:ladder} would give about $0.76H'$ rungs, and rungs $0$ and $1$ (infinite divisibility and self decomposability of $T_t$ for every $t\in[0,0.2]$) would need only an explicit zero density bound for $H_t$ and a small computation.
\end{remark}

\begin{problem}\label{mw:prob:khinchin}
Is there an analogue of Theorem~\ref{mw:thm:khinchin} at $\sigma=\tfrac12$, that is, a factorisation of the tilted law of $T$ or of the P\'olya variable through a law built from the primes, with the explicit formula replacing the Euler product?
\end{problem}

\section{Numerical checks}\label{mw:sec:checks}

The numerical results of this chapter are recomputed by the programs \texttt{maxwald\_*.py} of the repository (\texttt{code/\allowbreak README.md}). The zero counts of Table~\ref{mw:tab:Tstar} are the most delicate, and the other checks use multiprecision arithmetic, except the grids of Figure~\ref{mw:fig:offline}, the grid values and the gamma factors of the finite stages, and the sums and integrals of Table~\ref{mw:tab:AX}, which are in double precision; those of Table~\ref{mw:tab:AX} are checked against the sum over the zeros, to $5\times10^{-7}$. Table~\ref{mw:tab:checks} lists the checks of the results kept in this chapter.

\begin{table}[htbp]
\centering\small
\fitwidth{\begin{tabular}{p{2.1in}p{1.25in}p{1.3in}}
\toprule
result & test points & agreement \\
\midrule
kernel identity of Proposition~\ref{mw:prop:sum} & six values of $u$ & 24 digits \\
$\E M_4^{s/2}=2(1-2^{1-s})\xi(s)/(s-1)$ & $s=0.5,\,2,\,3+2i$ & 8 to 30 digits \\
Theorem~\ref{mw:thm:khinchin}, Mellin form & $\sigma=1.5,\,2,\,3.5$, $t=1,\,5,\,20$ & 30 digits \\
low coefficients of $\mathcal P_N$ equal those of $\vartheta_4$ & $N=4,8,16$ & exact \\
zero of $\xi_{64}$ near $\rho_1$ & $N=64$ & $|\re-\tfrac12|=5.2\times10^{-7}$ \\
zero counts of $F^{\mathrm{sym}}_N$ in $[-1,2]\times[0,T]$ against sign changes on the line (Table~\ref{mw:tab:Tstar}) & $N=8$, $16$, $32$, $48$, $64$, $96$, $128$ & windows of height $10$ \\
unit mass of the term of one zero, Proposition~\ref{mw:prop:lambda}(b) & $\alpha=0.55$ & 10 digits \\
width at half height of the term of $\rho_1$, Proposition~\ref{mw:prop:lambda} & $\alpha=0.55$ & 6 digits, up to the factor $1+\delta^2/(2\gamma_1^2)$ \\
$\gamma_1^2$ by Laplace inversion, Proposition~\ref{mw:prop:gnedenko} & $t\in[0.04,0.08]$ & $199.790$ \\
\bottomrule
\end{tabular}}
\caption{Numerical verification of results of this chapter that are not covered by the book's tables.}
\label{mw:tab:checks}
\end{table}
